\section{A time-weighted literature corollary and its order}
\label{sec:time-weighted}

\begin{corollary}[Time-weighted decay from Erd\'elyi]
\label{lem:ub-delay}
Let $(E,\nu)$ be a measure space, $N\ge1$, and
$\zeta_j\in\mathbb C$ with $\operatorname{Re}\zeta_j\ge1$.
For measurable finite coefficients define
\[
 f(t,z)=\sum_{j=1}^N c_j(z)e^{-\zeta_jt},\quad
 k(t)=\int_E|f(t,z)|^3\,d\nu(z),\quad K_T=\int_0^Tk(t)\,dt.
\]
If $T>0$ and $K_T<\infty$, then
\begin{equation}\label{eq:ub-delay}
 \int_0^Tt\,k(t)^{2/3}\,dt\le C_*N^{4/3}K_T^{2/3},
 \qquad C_*=4^{-1/3}36^{4/3}+1.
\end{equation}
No upper bound on the real parts or bound on the imaginary parts is needed.
\end{corollary}

\begin{proof}
Merge coincident exponents; estimates with the original $N$ follow by
monotonicity. Erd\'elyi's infinite--finite range theorem
\cite[Theorem 3.1, arXiv v1]{Erdelyi2017} applies to complex exponential
sums with nonpositive real parts. Applying it to $e^tf(t,z)$ gives
\[
 J(z):=\int_0^\infty e^t|f(t,z)|^2\,dt
 \le8191\int_0^{9N}e^t|f(t,z)|^2\,dt
 \le8191e^{9N}I_{9N}(z),
\]
where $I_b(z)=\int_0^b|f(t,z)|^2\,dt$. These integrals are finite
pointwise because the coefficients are finite and the real parts of the
decay rates are at least one. The midpoint estimate
\cite[Lemma 12.5, arXiv v1]{Erdelyi2017}, on an interval of length two,
applied to $e^{t/2}f(t,z)$, gives for $t\ge1$
\[
 e^t|f(t,z)|^2\le N\int_{t-1}^{t+1}e^u|f(u,z)|^2\,du\le NJ(z).
\]
Thus $|f(t,z)|^2\le8191Ne^{9N-t}I_{9N}(z)$.
This midpoint estimate allows arbitrary complex exponents.

Take $S=36N$. If $T\ge S$, Minkowski and H\"older give
\[
 \|I_{9N}\|_{L^{3/2}(E)}
 \le\int_0^{9N}k(t)^{2/3}\,dt
 \le(9N)^{1/3}K_T^{2/3}.
\]
The integral over $[S,T]$ is therefore at most
\[
 8191N(9N)^{1/3}(36N+1)e^{-27N}K_T^{2/3}
 \le N^{4/3}K_T^{2/3}.
\]
For the last inequality, $(36N+1)e^{-27N}$ decreases for $N\ge1$,
and $8191\,9^{1/3}37e^{-27}<1$.
For either $T\ge S$ or $T<S$, the initial part is bounded by
\[
 \int_0^{\min(T,S)}t k(t)^{2/3}\,dt
 \le\left(\int_0^S t^3\,dt\right)^{1/3}K_T^{2/3}
 =4^{-1/3}(36N)^{4/3}K_T^{2/3}.
\]
This proves the claim, including $K_T=0$.
\end{proof}

\begin{proposition}[An explicit order witness]\label{prop:time-order}
For each $N\ge1$ there is a real sum with exactly $N$ nonzero terms,
all decay rates in $[1,5/4]$, such that, on the single-point space and
with $T=N$,
\begin{equation}\label{eq:time-order}
 \frac{\int_0^N t f_N(t)^2\,dt}
 {\left(\int_0^N f_N(t)^3\,dt\right)^{2/3}}
 \ge\frac{N^{4/3}}{128}.
\end{equation}
Consequently the order $N^{4/3}$ in \eqref{eq:ub-delay} cannot be
replaced by $o(N^{4/3})$ in its abstract class.
\end{proposition}

\begin{proof}
Set $\epsilon=N^{-2}$, $w(t)=(1-e^{-\epsilon t})/\epsilon$, and
\[
 f_N(t)=e^{-t}\sum_{k=0}^{N-1}\frac{w(t)^k}{k!}
       =\sum_{j=0}^{N-1}b_je^{-(1+j\epsilon)t},\qquad
 b_j=\frac{(-1)^j}{j!\epsilon^j}
      \sum_{h=0}^{N-1-j}\frac1{\epsilon^h h!}.
\]
The coefficient identity follows by the binomial theorem and all its
inner summands are positive. Thus every $b_j\ne0$, and the distinct
decay rates lie in $[1,5/4]$ since $(N-1)/N^2\le1/4$.
Write $f_{0,N}(u)=e^{-u}\sum_{k=0}^{N-1}u^k/k!$.
This is the probability that a Poisson variable of mean $u$ is less
than $N$, equivalently the elementary incomplete-gamma identity
\cite[\S8.4.11]{DLMF}. For $0\le t\le N/2$ one has
\[
 0\le w(t)\le t,\qquad t-w(t)\le\epsilon t^2/2\le1/8,
 \qquad f_{0,N}(w(t))\ge1-w(t)/N\ge\tfrac12.
\]
The middle inequality follows by integrating
$1-e^{-\epsilon u}\le\epsilon u$; the last one is Markov's probability
inequality. Hence
$f_N(t)=e^{-(t-w(t))}f_{0,N}(w(t))\ge\tfrac12e^{-1/8}>1/4$ there.
On all $t\ge0$, $0<f_N(t)\le f_{0,N}(t)\le1$.
Therefore $K_N=\int_0^N f_N^3\,dt$ lies in $(0,N]$ and
\[
 \int_0^N t f_N^2\,dt\ge\frac1{16}\int_0^{N/2}t\,dt
 =\frac{N^2}{128}.
\]
Division proves \eqref{eq:time-order}. Nothing excludes $N=1$:
then $b_0=1$ and $f_1=e^{-t}$. If unit cubic mass is desired,
replace $f_N$ by $K_N^{-1/3}f_N$; the quotient is unchanged.
\end{proof}

This elementary construction concerns the abstract temporal norm quotient.
Its nearby real decay rates need not be the squared integer frequencies
of a shear, and it has no corresponding spatial constraint. It does not
prove optimality of the $q^7$ flux coefficient, a sharp exponential
constant, or a dynamical transition. The upper estimate is a direct
literature corollary; no priority claim is made for this order witness.

For later application, if the temporal rates of $G(s,z)$ have real
parts at least $\lambda>0$, write
$h(s)=\int_E|G(s,z)|^3\,d\nu$, $H=\int_0^4h(s)\,ds$.
For $0\le\eta(s)\le C_\eta s$, substitution $t=\lambda s$ gives
\begin{equation}\label{eq:ub-cutoff-time}
 \lambda\int_0^4\eta(s)h(s)^{2/3}\,ds
 \le C_* C_\eta N^{4/3}\lambda^{-1/3}H^{2/3}.
\end{equation}
